% ============================================================
%  Apêndices
% ============================================================

\chapter[Apêndices]{Apêndices}

\section*{Apêndice A: Cronologia dos Desenvolvimentos Escatológicos}
\addcontentsline{toc}{section}{Apêndice A: Cronologia dos Desenvolvimentos Escatológicos}

\begin{description}[leftmargin=3.2cm,style=multiline]
  \item[40.000--10.000 a.C.] Paleolítico Superior: primeiras práticas funerárias intencionais.
  \item[c. 9.500 a.C.] Göbekli Tepe: iconografia ritual associada à morte e aos ancestrais.
  \item[3.500--3.000 a.C.] Primeiros textos cuneiformes sobre morte e submundo.
  \item[2.700--2.500 a.C.] Tradições de Gilgamesh: descrições literárias de descidas ao submundo.
  \item[2.100--2.000 a.C.] \emph{Descida de Inanna}: sistematização da geografia do mundo inferior.
  \item[c. 1.750 a.C.] Código de Hammurabi: precedentes de retribuição proporcional.
  \item[1.550--1.077 a.C.] Livro dos Mortos egípcio: julgamento moral explícito.
  \item[1.200--800 a.C.] Tradições homéricas: consolidação literária do Hades.
  \item[586--538 a.C.] Exílio babilônico e transformações da escatologia judaica.
  \item[Séc. VI a.C.] Mistérios órficos: julgamento moral individual.
  \item[332--63 a.C.] Período helenístico: florescimento da literatura apocalíptica judaica.
  \item[30--100 d.C.] Novo Testamento: formulações escatológicas cristãs primitivas.
  \item[325] Concílio de Niceia: consolidação de fundamentos cristológicos.
  \item[451] Concílio de Calcedônia: definições cristológicas posteriores.
  \item[553] Segundo Concílio de Constantinopla e controvérsias sobre o origenismo.
  \item[711--1492] Al-Andalus: intercâmbios intelectuais islâmico-cristãos.
  \item[1054] Grande Cisma: aprofundamento das divergências entre Oriente e Ocidente.
  \item[c. 1308--1320] \emph{Divina Comédia}: síntese medieval do imaginário infernal.
  \item[1517] Reforma protestante e contestação da doutrina do purgatório.
  \item[1545--1563] Concílio de Trento: reafirmação da doutrina católica.
  \item[1859] \emph{A origem das espécies}: novos desafios às cosmologias tradicionais.
  \item[1900] \emph{A interpretação dos sonhos}: expansão das leituras psicológicas.
  \item[1945] Holocausto: transformação moderna da imaginação do mal.
  \item[1962--1965] Concílio Vaticano II: renovação da linguagem teológica católica.
  \item[Década de 1990] Popularização da internet e digitalização do imaginário escatológico.
\end{description}

\section*{Apêndice B: Comparação de Tradições Escatológicas}
\addcontentsline{toc}{section}{Apêndice B: Comparação de Tradições Escatológicas}

\small
\renewcommand{\arraystretch}{1.2}
\begin{longtable}{@{}p{2.4cm}p{3.1cm}p{2.6cm}p{3.0cm}p{3.2cm}@{}}
  \toprule
  \textbf{Tradição} & \textbf{Localização} & \textbf{Governante} & \textbf{Critério de acesso} & \textbf{Destino ou punição} \\
  \midrule
  \endfirsthead
  \toprule
  \textbf{Tradição} & \textbf{Localização} & \textbf{Governante} & \textbf{Critério de acesso} & \textbf{Destino ou punição} \\
  \midrule
  \endhead
  Mesopotâmica & Irkalla subterrâneo & Ereshkigal & Condição universal dos mortos & Existência diminuída \\
  Egípcia & Duat & Osíris & Julgamento moral & Bem-aventurança ou segunda morte \\
  Greco-romana & Hades e Tártaro & Hades/Plutão & Conduta e estatuto heroico & Sombra comum ou punição exemplar \\
  Nórdica & Helheim & Hel & Tipo de morte & Existência sombria; guerreiros podem alcançar outros salões \\
  Maia & Xibalba & Senhores de Xibalba & Provas e conhecimento ritual & Derrota, transformação ou transcendência \\
  Asteca & Mictlan & Mictlantecuhtli & Tipo de morte e travessia ritual & Repouso após jornada pelo submundo \\
  Judaica & Sheol/Geena & Deus como juiz & Aliança, justiça e julgamento & Morte comum, purificação ou punição \\
  Cristã & Inferno e purgatório & Deus; figuras demoníacas subordinadas & Fé, graça e conduta moral & Punição, purificação ou separação de Deus \\
  Islâmica & Jahannam & Deus como juiz & Fé, obras e misericórdia divina & Punição graduada, segundo diferentes leituras \\
  \bottomrule
\end{longtable}
\normalsize
